\section{Step 15 Record the discovery predecessor and distance}\label{maze-s15}

Discovery assigns the candidate to the next BFS layer. Its parent is the cell being processed, and its distance is one more than that parent distance.

\subsection{Write or inspect}

Append the state, parent, and distance stores. Keep the mark before the enqueue operation.

\MARSCodeLabel{MIPS assembly}
\begin{minted}{mars-idea}
    li $t1, 1
    sb $t1, 0($t0)
    daddu $t0, $s2, $t3
    sb $s7, 0($t0)
    daddu $t0, $s3, $s7
    lbu $t1, 0($t0)
    addiu $t1, $t1, 1
    daddu $t0, $s3, $t3
    sb $t1, 0($t0)
\end{minted}

\subsection{Expected observation}

A candidate that passes the checks will be marked waiting and receive a predecessor and distance. With the default start, the first north candidate is rejected, so these stores are not yet observed. Step 17 completes the direction loop and reaches the other neighbors.

\subsection{Check your understanding}

Why mark the cell before a later dequeue?

\begin{itemize}
\item \textbf{A} It makes all edges cost zero.
\item \textbf{B} It removes every wall.
\item \textbf{C} It prevents duplicate enqueues through other incoming edges.
\end{itemize}
